% ATTEMPT_STATUS: unresolved
% ATTEMPT_ORIGIN: new research, 2026-10-01
% TOP_PROBLEM: {"id": 1402, "title": "Albertson Conjecture", "category_id": 3, "category": {"id": 3, "name": "graph_theory", "display_name": "Graph Theory", "description": "Problems involving graphs, networks, and their properties.", "slug": "graph-theory", "order_index": 3, "created_at": "2026-07-31T15:26:25.671Z"}, "status": "open"}
% FIRST_POSTED: 2026-10-01
% Imported from: unsolvedmath_additional_500.tex; UnsolvedMath ID 1402
% !TeX program = lualatex
% Compile twice: lualatex 1402.tex
% Self-contained: no external bibliography, images, or \input files.
% Numeric section labels and visible numbers are original UnsolvedMath IDs.
\documentclass[11pt,a4paper]{article}
\usepackage[margin=24mm,headheight=14pt]{geometry}
\usepackage{fontspec}
\setmainfont{Latin Modern Roman}
\setsansfont{Latin Modern Sans}
\setmonofont{Latin Modern Mono}
\usepackage{amsmath,amssymb,mathtools}
\usepackage{unicode-math}
\setmathfont{Latin Modern Math}
\usepackage{microtype}
\usepackage{enumitem,xurl,needspace}
\usepackage[dvipsnames]{xcolor}
\usepackage{hyperref}
\hypersetup{unicode=true,colorlinks=true,linkcolor=MidnightBlue,urlcolor=MidnightBlue,
 pdftitle={UnsolvedMath Problem 1402},pdfauthor={Source-annotated compilation},bookmarksnumbered=true}
\usepackage{bookmark,tocloft,titlesec,fancyhdr}
\setlength{\cftsecnumwidth}{4.2em}
\cftsetpnumwidth{2.5em}
\cftsetrmarg{4em}
\titleformat{\section}{\Large\bfseries\raggedright}{\thesection}{0.7em}{}
\pagestyle{fancy}\fancyhf{}
\fancyhead[L]{\small\sffamily UnsolvedMath research attempt}
\fancyhead[R]{\small Original catalogue IDs}
\fancyfoot[C]{\thepage}
\setlength{\parindent}{0pt}
\setlength{\parskip}{5pt plus 1pt}
\setlength{\emergencystretch}{3em}
\setcounter{tocdepth}{1}\setcounter{secnumdepth}{1}
\setlist[itemize]{leftmargin=3em,itemsep=4pt,topsep=4pt,parsep=0pt}
\allowdisplaybreaks
\newcommand{\N}{\mathbb{N}}
\newcommand{\Z}{\mathbb{Z}}
\newcommand{\Q}{\mathbb{Q}}
\newcommand{\R}{\mathbb{R}}
\newcommand{\C}{\mathbb{C}}
\newcommand{\F}{\mathbb{F}}
\newcommand{\A}{\mathbb{A}}
\newcommand{\PP}{\mathbb{P}}
\newcommand{\E}{\mathbb{E}}
\newcommand{\eps}{\varepsilon}
\newcommand{\dd}{\,\mathrm d}
\newcommand{\sref}[2]{\hyperlink{source:#1:#2}{[S#2]}}
\newif\ifshowcatalogue
\showcataloguetrue
\newcommand{\cataloguescope}[1]{\ifshowcatalogue
 \paragraph{Statement recorded in the supplied source.}
 #1\fi}
\begin{document}
% UnsolvedMath ID 1402; source catalogue code GRAPH-015

\setcounter{section}{1401}

\section{Albertson’s Crossing-Number Conjecture}\label{1402}

{\small\sffamily UnsolvedMath ID 1402 · Graph Theory}

\subsection{Definitions and mathematical statement}

\paragraph{Shared notation.}Unless a section says otherwise, $\N=\{1,2,\ldots\}$,
$\Z$ is the ring of integers, $\Q,\R,\C$ are the rational, real, and complex
fields, $[N]=\{1,\ldots,\lfloor N\rfloor\}$, and $\log$ is the natural
logarithm. For real functions of a parameter tending to infinity,
$f\ll g$ and $f=O(g)$ mean $|f|\leq Cg$ eventually for some fixed $C$;
$f\gg g$ reverses this comparison; $f=o(g)$ means $f/g\to0$;
$f\sim g$ means $f/g\to1$; and $f\asymp g$ means both order bounds.
A subscript on $O$ or $\ll$ specifies allowed constant dependence.
The expression $n^{o(1)}$ in an upper bound means growth smaller than
$n^\epsilon$ for each fixed $\epsilon>0$. Local definitions govern any
exception, finite-field convention, or use of the same symbol for another
object. An asymptotic claim is not automatically an effective algorithm.



The crossing number $\operatorname{cr}(G)$ is the least number of proper edge crossings over plane drawings with vertices distinct and edges drawn as arcs, in general position. The conjecture is $\operatorname{cr}(G)\geq\operatorname{cr}(K_{\chi(G)})$ for every finite graph $G$. Straight-line drawings are not required. \sref{1402}{1} \sref{1402}{2}

\paragraph{Statement recorded in the supplied source.}
Can the crossing number of a graph be lower-bounded by the crossing number of a complete graph with the same chromatic number? \sref{1402}{1}

\paragraph{Residual task reported by the archive.}

The embedded review (2026-08-17) records: Establish the crossing-number lower bound for every chromatic number and graph order. \sref{1402}{1}

\subsection{Short English statement}

Must every graph incur at least as many unavoidable plane crossings as a complete graph requiring the same number of vertex colours?

\subsection{Research attempt}
\textbf{Status: UNRESOLVED.} We reduce a possible counterexample to a vertex-critical graph, derive an elementary crossing lemma, and prove the conjectured comparison when that critical graph has sufficiently many vertices relative to its chromatic number. This leaves a concrete finite range of orders for each fixed chromatic number, but does not eliminate it.

\paragraph{Literature check.}
Cranston [S2], primary abstract consulted 1 October 2026, reports the conjecture through chromatic number 24 and stronger restrictions on the order of a minimum counterexample. We do not reproduce those results. The calculation below is a weaker self-contained reduction illustrating the role of criticality and crossing density. No improvement of the published constants is claimed.

\paragraph{Passing to a critical subgraph.}
Let $r=\chi(G)$. Repeatedly delete a vertex whose deletion leaves chromatic number $r$, stopping at a graph $H$ for which every vertex deletion lowers it. Then $\chi(H)=r$, and $\operatorname{cr}(H)\le\operatorname{cr}(G)$ by restricting a drawing. Every vertex of $H$ has degree at least $r-1$: otherwise an $(r-1)$-colouring of $H-v$ would leave a colour available for $v$. Writing $n=|H|$ and $m=|E(H)|$ gives $m\ge(r-1)n/2$.

If $n=r$, the graph must be complete. Indeed any missing edge permits its endpoints to share a colour while every other vertex receives its own, giving at most $r-1$ colours. Thus order $r$ itself satisfies the desired comparison exactly. The possible difficulty starts at larger critical graphs.

\paragraph{An elementary crossing-number inequality.}
We may take a drawing realizing the crossing number in which no two adjacent edges cross and every crossing is between two edges with four distinct endpoints. A crossing between adjacent edges can be removed by exchanging their initial arcs and making a local perturbation, without increasing other intersections. Degenerate multiple crossings can also be separated into pairwise crossings.

In any such drawing, removing at most one edge for each crossing leaves a planar simple graph. The planar edge bound, weakened to $3n$ so that all small orders are included, therefore yields $c\ge m-3n$, where $c$ is the number of crossings. Now retain each vertex independently with probability $p$ and restrict the drawing to the induced subgraph. Expected vertex, edge, and crossing counts are $pn,p^2m,p^4c$. Applying the same weak planar inequality to each sampled drawing and averaging gives
\[
p^4c\ge p^2m-3pn.
\]
For $m\ge4n$, choose $p=4n/m\le1$. Substitution gives
\[
\operatorname{cr}(H)\ge\frac{m^3}{64n^2}.
\]
The fourth power of $p$ is why the four-distinct-endpoints convention had to be justified. Using $p^4$ for a crossing of adjacent edges would be incorrect.

\paragraph{A range where the target follows.}
For $r\ge9$, criticality ensures $m\ge4n$, hence
\[
\operatorname{cr}(H)\ge\frac{(r-1)^3n}{512}.
\]
Place $r$ points in strictly convex position and draw all straight chords. Every four-set determines exactly one crossing between its diagonals, so this gives a drawing of $K_r$ with $\binom r4$ crossings. Thus
\[
\operatorname{cr}(K_r)\le\binom r4.
\]
It follows that the desired inequality holds whenever
\[
n\ge\frac{512\binom r4}{(r-1)^3}.
\]
This sufficient threshold is approximately $(64/3)r$ for large $r$. The coefficient is intentionally weak; the source reports much sharper order reductions. The deduction does not require knowing the exact crossing number of $K_r$, whose general formula is itself a separate unresolved target.

\paragraph{A second, weaker inequality for moderate order.}
The unsampled planar inequality combined with criticality gives
\[
\operatorname{cr}(H)\ge \frac{(r-7)n}{2}.
\]
For $r>7$ this is positive, but grows only linearly with $r$ at fixed $n/r$. It generally cannot match a quartic complete-graph crossing scale. This makes precise why simply deleting one edge per crossing and applying a colouring degree bound is insufficient in the critical order range.

For $r\le4$, $K_r$ has a plane drawing, so $\operatorname{cr}(K_r)=0$ and the target follows trivially. No appeal to the converse assertion that every planar graph is four-colourable is needed for this direction.

\paragraph{Verification and remaining gap.}
The random sampling is used only as an averaging proof; no simulation is claimed. The parameter choice has been restricted to $m\ge4n$ so that it is a probability. Monotonicity is used for taking subgraphs, where it is immediate by deleting edges and vertices from a drawing; no unsupported claim about arbitrary graph contractions is used.

For each fixed $r\ge9$, the remaining candidates after this elementary argument have critical order strictly between $r$ and the displayed threshold. Finiteness at fixed $r$ does not give one finite computation for all $r$. The missing step is either a stronger structural lower bound for crossings in that range or a verified counterexample drawing. Our work reduces the search but does not settle Albertson's conjecture.

\subsection{Sources}

{\small

\begin{itemize}

\item[\hypertarget{source:1402:1}{S1}] ULAM.AI, \emph{UnsolvedMath}, supplied \texttt{problems.json}, record 1402 (GRAPH-015), statement and background. The embedded literature-review date is 2026-08-17; that is the archive’s review date, not a fresh certification. Dataset landing page: \url{https://huggingface.co/datasets/ulamai/UnsolvedMath}.

\item[\hypertarget{source:1402:2}{S2}] D. W. Cranston, Progress on Albertson's Conjecture, arXiv:2512.08020 (2025). \url{https://arxiv.org/abs/2512.08020}. Bibliographic information imported from the supplied record.

\item[\hypertarget{source:1402:3}{S3}] J. Fox, J. Pach and A. Suk, Immersions and Albertson's Conjecture, SoCG 2025, \nolinkurl{doi:10.4230/LIPIcs.SoCG.2025.50.} \url{https://drops.dagstuhl.de/entities/document/10.4230/LIPIcs.SoCG.2025.50}. Bibliographic information imported from the supplied record.

\end{itemize}

}

\label{end:1402}
\end{document}
